\documentclass[a4paper,12pt]{article}

\usepackage[utf8]{inputenc}
\usepackage[T1]{fontenc}
\usepackage[polish]{babel}
\usepackage{amsmath, amssymb}
\usepackage{graphicx}
\usepackage{float}
\usepackage{hyperref}
\usepackage{geometry}
\usepackage{csvsimple} 

\geometry{margin=2.5cm}

\title{Sprawozdanie 1 z Statystyki i Modeli Liniowych}
\author{Wiktor Zoga}
\date{\today}

\begin{document}

\maketitle
\newpage

\section{Zadanie 4}

\subsection{Opis eksperymentu}
W tym zadaniu w zależności od wartości parametrów: 

\begin{itemize}
    \item $\mu=0, \theta=1$
    \item $\mu=0, \theta=4$
    \item $\mu=4, \theta=1$
    \item $\mu=4, \theta=4$
\end{itemize}

esytmowano wariancję próby wielkości $n=50$ z rozkładu $N(\mu, \theta)$,

przy pomocy trzech esymatorów:

\begin{itemize}
    \item $\hat{\theta}_1 = S^{2} = \frac{1}{n-1}\sum_{i=1}^{n}(X_{i}-\bar{X})^{2}$
    \item $\hat{\theta}_2 = \frac{1}{n-1}\sum_{i=1}^{n}(X_{i} - median(X_{1}, \dots, X_{n}))^{2}$
    \item $\hat{\theta}_3 = median((X_1 - \bar{X})^{2}, \dots, (X_{n} -\bar{X})^2)$
\end{itemize}

Doświadczenie powtórzono $10{,}000$ razy. Porównano esymatory przy pomocy wyrkesów pudełkowych oraz
wyliczono wariancję, błąd średniokwadratowy oraz obciązenie każdego z esymatorów. 

\section{Wyniki}

\newcommand{\plotwithstats}[2]{
    \subsection{$\mu=#1$, $\theta=#2$}
    \begin{figure}[H]
        \centering
        \includegraphics[width=0.9\textwidth]{../plots/boxplot_mu_#1_theta_#2.png}
        \caption{Boxplot estymatorów dla $\mu=#1$, $\theta=#2$.}
    \end{figure}

    \begin{center}
    \csvreader[
        tabular=|l|c|c|c|c|,
        table head=\hline Estymator & Średnia & Wariancja & Bias & MSE \\ \hline,
        late after line=\\\hline,
        filter expr={true},  % brak filtrowania
        respect percent
    ]{../plots/stats_mu_#1_theta_#2.csv}{Estimator=\est,Mean=\mean,Variance=\var,Bias=\bias,MSE=\mse}
    {\est & \pgfmathprintnumber{\mean} & \pgfmathprintnumber{\var} & \pgfmathprintnumber{\bias} & \pgfmathprintnumber{\mse} }
    \end{center}

    \newpage
}

\plotwithstats{0}{1}
\plotwithstats{0}{4}
\plotwithstats{4}{1}
\plotwithstats{4}{4}

\section{Dyskusja Wyników}

W tym zadaniu celem była estymacja parametru $\theta$, 
który w notacji rozkładu $\mathcal{N}(\mu, \theta)$ odpowiada \textbf{wariancji}.

\subsection*{Analiza Estymatora $\hat{\theta}_1 = S^2$}

Jest to klasyczny, \textbf{nieobciążony estymator wariancji} z próby.
\begin{itemize}
    \item \textbf{Obciążenie:} Zgodnie z rachunkami, które przeprowadziliśmy na wykładzie $E[S^2] = E[\hat{\theta}_1] = \sigma^2 = \theta$. 

    \item \textbf{MSE:} Ponieważ estymator jest nieobciążony, jego MSE jest równe jego wariancji: $\text{MSE}(\hat{\theta}_1) = \text{Var}(\hat{\theta}_1)$.

    Dokładnie to:
    \[
        \boxed{
            E[S^2] = \sigma^2, \quad \mathrm{Var}(S^2) = \frac{2\sigma^4}{n-1}
        }       
    \]
    \begin{itemize}
        \item Gdy $\theta = 1$ (wariancja populacji), $\text{MSE}(\hat{\theta}_1) = \frac{2 \cdot 1^2}{49} \approx 0.04$.
        \item Gdy $\theta = 4$ (wariancja populacji), $\text{MSE}(\hat{\theta}_1) = \frac{2 \cdot 4^2}{49} = \frac{32}{49} \approx 0.65$.
    \end{itemize}
    Wyniki eksperymentu zgadzają się z teoretycznymi obliczeniami.

\end{itemize}

\subsection*{Analiza Estymatora $\hat{\theta}_2$}

Estymator ten jest analogiczną konstrukcją do $S^2$, jednak zamiast optymalnego estymatora średniej $\bar{X}$, używa on estymatora odpornego -- \textbf{mediany} ($\hat{m} = \text{median}(X)$).
\begin{itemize}
    \item W rozkładzie normalnym $\hat{m}$ jest również nieobciążonym estymatorem $\mu$, jednak jest on \textbf{mniej efektywny} (ma większą wariancję) niż $\bar{X}$.
    \item Ponieważ do obliczenia $\hat{\theta}_2$ używamy "mniej dokładnego" oszacowania środka rozkładu, intuicyjnie spodziewamy się, że $\hat{\theta}_2$ będzie \textbf{nieco gorszym} estymatorem wariancji $\theta$ niż $\hat{\theta}_1$.
    \item Wyniki symulacji sugerują, że jest to estymator poprawny i bardzo zbliżony do $\hat{\theta}_1$, jednak różnica w biasie i MSE jest zauważalna.
\end{itemize}

\subsection*{Analiza Estymatora $\hat{\theta}_3$}

Estymator $\hat{\theta}_3 = \text{median}((X_i - \bar{X})^2)$ jest przykładem estymatora \textbf{obciążonego}.

\subsubsection*{Analiza asymptotyczna}

Aby formalnie wykazać obciążenie, zbadajmy granicę, do której zbiega estymator. Wraz ze wzrostem $n \to \infty$, średnia z próby $\bar{X}$ zbiega do $\mu$. Zatem $\hat{\theta}_3$ zbiega (według prawdopodobieństwa) do $m = \text{median}((X_i - \mu)^2)$.

Musimy wyznaczyć wartość tej mediany $m$.
\begin{itemize}
    \item Niech $X_i \sim \mathcal{N}(\mu, \theta)$.
    \item Standaryzujmy zmienną: $Z = (X_i - \mu) / \sqrt{\theta}$, gdzie $Z \sim \mathcal{N}(0, 1)$.
    \item Kwadrat tej zmiennej, $Z^2$, ma rozkład $\chi^2_1$ (Chi-kwadrat z 1 stopniem swobody).
\end{itemize}
Wartość $m$, do której zbiega estymator, to:
\[ m = \text{median}((X_i - \mu)^2) = \text{median}(\theta \cdot Z^2) = \theta \cdot \text{median}(\chi^2_1) \]

Mediana rozkładu $\chi^2_1$ to wartość $t$, dla której $P(Z^2 \le t) = 2\Phi(\sqrt{t}) - 1 = 0.5$. Jak wynika z obliczeń, $t = (\Phi^{-1}(0.75))^2 \approx (0.6745)^2 \approx \mathbf{0.4549}$.

\subsubsection*{Wnioski z symulacji}

Oznacza to, że estymator $\hat{\theta}_3$ nie zbiega do $\theta$, lecz do $\approx 0.4549 \cdot \theta$.
Fakt ten w pełni wyjaśnia wyniki obserwowane w symulacji:
\begin{itemize}
    \item Estymator wykazuje \textbf{znaczące obciążenie}, systematycznie zaniżając prawdziwą wariancję $\theta$ o około 55\%.
    \item Obserwacja z eksperymentu dla $\theta=4$ (gdzie wynik był $\approx 1.86$) jest bardzo bliska wartości teoretycznej, $4 \cdot 0.4549 \approx 1.82$.
    \item To duże obciążenie jest głównym składnikiem \textbf{wysokiego błędu średniokwadratowego (MSE)}, zgodnie ze wzorem $\text{MSE} = \text{Var} + \text{Bias}^2$.
\end{itemize}

\end{document}
